
\documentclass[12pt,norsk,a4paper]{article}
\usepackage[norsk]{babel}
\usepackage[norsk]{isodate}
\usepackage[T1]{fontenc}
\usepackage{parskip}
\usepackage{url}
\author{fridatbe}
\title{IN1080 --- Obligatorisk oppgave 2} % erstatt N
\date{\today}
\usepackage[small,euler-digits]{eulervm}
\usepackage{color}|
% \usepackage{minted}
\usepackage{xcolor} % to access the named colour LightGray
\definecolor{LightGray}{gray}{0.9}
\usepackage{a4wide}
\usepackage[top=2cm]{geometry}

% 6. Figurer og bilder

% 6a. For å bryte tekst rundt en figur
%\usepackage{wrapfig}

% 6b. For bilder - med denne pakken kan du legge inn bilder og
% illustrasjoner slik:
%                  \includegraphics[width=.5\textwidth]{bilde.pdf}
\usepackage{graphicx}
\graphicspath{ {./fig/} }

% 7. For grafer og slikt
%\usepackage{tikz}
%\usetikzlibrary{trees}
% Det fins mange tikz-bibliotek, det er nesten ingen grenser for hva
% du kan lage med dette. Se http://www.texample.net/tikz/

% 8. Egendefinerte kommandoer kan du legge inn slik:
\newcommand{\tuple}[1]{\ensuremath{\langle #1\rangle}}
\newcommand{\set}[1]{\ensuremath{\{#1\}}}
\newcommand{\imp}{\ensuremath{\rightarrow}}
\newcommand{\M}{\ensuremath{\mathcal{M}}}
\newcommand{\AF}{edge from parent[draw=none]}
\newcommand{\NODE}[1]{\{node \{\ensuremath{#1}\}\}}

% 9. Egendefinert overskrift som passer obligformatet
\usepackage[explicit]{titlesec}

\titleformat{\section}{\normalfont\Large\bfseries}{}{0em}{#1\ \thesection}

% http://en.wikibooks.org/wiki/LaTeX/Customizing_LaTeX#New_commands

%%%%%%%%%%%%%%%%%%%%%%%%%%%%%%%%%%%%%%%%%%%%%%%%%%%%%%%%%%%%%%%%%%%%%%%%%%%
% Selve innholdet:
\begin{document}

    % a. Lager en tittel på dokumentet.
    \maketitle

    % a. For å få bokstavnummerering på 'subsection's
    \renewcommand{\thesubsection}{(\alph{subsection})}

    \section{Oppgave} % for hver oppgave, lag en \section{Oppgave}

    Tekst kan skrives inn direkte. Du trenger ikke   ta hensyn  til
    hvor mye luft det   er   mellom ordene.

    \subsection*{Formler}

    Formler kan du blant skrive inn slik: $V=R*I$.

    Fremheving av formler kan du gjøre som i \ref{eq:kirchoff}:

    \begin{equation}
        \sum_{k=0}^n V_k = 0
        \label{eq:kirchoff}
    \end{equation}

    Legg merke
    til at du må ha et mellomrom etter kommandoen - dette forsvinner ved
    kompilering. Tvungent mellomrom skrives slik: \ .


    \subsection*{Lister}

    \begin{description}
        \item[(a)] Det er flere måter å lage lister på.
        \item[(b)] Denne listen er laget med {\tt description}.
        \item[(c)] Du kan også bruke {\tt itemize} og {\tt enumerate}.
    \end{description}

    \subsection*{Tabeller}

    For tabeller kan du bruke {\tt tabular}. Eksempel på en
    sannhetsverditabell:

    \begin{tabular}{ccc}
        $P$ & $Q$ & $P\land Q$ \\
        \hline
        $1$ & $1$ & ${\color{red}1}$ \\
        $1$ & $0$ & ${\color{red}0}$ \\
        $0$ & $1$ & ${\color{red}0}$ \\
        $0$ & $0$ & ${\color{red}0}$ \\
    \end{tabular}

    \subsection*{Kode}

    \begin{minted}
        [
        frame=lines,
        framesep=2mm,
        baselinestretch=1.2,
        bgcolor=LightGray,
        fontsize=\footnotesize,
        linenos
        ]
        {python}
        from matplotlib import pyplot as plt

        U=[0, 1, 2, 3]
        I=[0, 0.02, 0.04, 0.06]

        plt.xlabel("Voltage (V)")
        plt.ylabel("Current (A)")
        plt.plot(U,I)
        plt.scatter(U,I)
        plt.grid()
        plt.savefig('./fig/plotPy.pdf')
    \end{minted}

    \subsection*{Bilder og figurer}

    \begin{figure}[h]
        \caption{I/V plot}
        \centering
        \includegraphics[width=.7\textwidth]{plotPy.pdf}
        \label{fig:plot}
    \end{figure}

    Hvis det er tekst som beskriver en figur eller et bilde er det god skikk å referere til bilde. For eksempel:  ``I figur \ref{fig:plot} ser vi [\ldots].''

\end{document}
